\documentclass[10pt,a4paper]{amsart}
\usepackage[margin=19mm]{geometry}
\usepackage{amsmath,amssymb,mathtools}
\usepackage[scr=rsfs,cal=euler]{mathalfa}
\usepackage{xfrac}
\usepackage[unicode,hidelinks,bookmarks=false]{hyperref}
\newcommand{\ie}{i.e.\ }
\newcommand{\cf}{cf.\ }
\newcommand{\resp}{respectively\ }
\newcommand{\Subsection}{\subsection}
\providecommand{\nobreakdash}{\nobreak\hbox{-}}
\setlength{\emergencystretch}{2em}
\multlinegap0pt
\usepackage{bm}
\usepackage{bbm}
\usepackage{dsfont}
\usepackage{xspace}
\usepackage{cancel}
\usepackage{mathtools}
\usepackage{cleveref}
\usepackage[shortlabels]{enumitem}
\usepackage{amsthm}
\makeatletter
\@ifundefined{thm}{\newtheorem{thm}{Thm}}{}
\@ifundefined{prop}{\newtheorem{prop}[thm]{Proposition}}{}
\makeatother
\newcommand{\fp}{{\mathfrak p}}
\newcommand{\rh}{{\rm h}}  % henselisation
\newcommand{\sh}{{\rm sh}}   % strict henselisation
\title[Tame fundamental groups of rigid spaces]{Tame fundamental groups of rigid spaces}
\author{Piotr Achinger, Katharina H\"ubner, Marcin Lara, Jakob Stix}
\date{Source: arXiv:2608.12601v1 (CC BY 4.0)}
\begin{document}
\maketitle

\noindent\emph{Source and licence.} Tame fundamental groups of rigid spaces. Version arXiv:2608.12601v1 (CC BY 4.0), distributed under the Creative Commons Attribution 4.0 International licence, as recorded in the arXiv licence metadata for this exact version and shown on the abstract page https://arxiv.org/abs/2608.12601v1. Full rights notice: licenses/arxiv-2608.12601-CC-BY-4.0.txt.\par\smallskip

\noindent\textit{Editorial scope.} Counted unit: the Prop whose source label is \texttt{composition-tame} in the article (source0.tex lines 853--878, source label \texttt{composition-tame}), delivered together with the article's own complete original proof of it (source0.tex lines 880--907, one \texttt{proof} environment of 28 lines). No mathematics is written, repaired or re-derived: statement and proof are the article's own text line for line. Required context retained uncounted: henselisation-composition (lines 829--841). Every definition, display and label of those lines is retained unchanged. Omitted from this package and not redistributed: 1--828, 842--852, 879--879, 908--4973. Nothing inside a retained excerpt is omitted or altered: every retained body is delivered exactly as the article's lines are written, so no omission receipt of any kind accompanies this package. Citations: the retained excerpts carry no citation command at all, so no bibliography entry of the article is transcribed and no reference is redistributed. Reference adaptation: none was needed and none was made; every reference inside the delivered unit resolves inside it, so no unresolved reference or citation is produced. Printed slips kept literal: no printed word, symbol, hypothesis or number of the counted statement or of its proof was altered. Layout and class: the article's own class is \texttt{amsart} with packages including mathrsfs, mathtools, amssymb, latexsym, eucal, extarrows, caption, dsfont, fullpage, graphicx, color, fontenc, inputenc, enumitem; the wrapper is the lane's pinned \texttt{amsart} editorial wrapper at 10pt on A4, which restates the theorem-like environments the retained text needs and adds the article's own macro definitions that the retained text needs (\texttt{\textbackslash fp}, \texttt{\textbackslash rh}, \texttt{\textbackslash sh}), with extra wrapper packages (bm, bbm, dsfont, xspace, cancel, mathtools, cleveref, enumitem). The delivered numbering is the unit's own. Assets: no retained span contains a figure environment, a graphics-inclusion command, a PDF-page inclusion, a table or a TikZ picture; no image file is copied into this package, the package declares no asset dependency and no image file is redistributed with it. Licence: version arXiv:2608.12601v1 of this article is distributed under the Creative Commons Attribution 4.0 International licence, as recorded in the arXiv licence metadata for this exact version and shown on the abstract page https://arxiv.org/abs/2608.12601v1. A full rights notice is delivered at licenses/arxiv-2608.12601-CC-BY-4.0.txt. Characters: every retained excerpt is pure ASCII, so the delivered engine, XeLaTeX with its default Latin Modern text font, reports no missing character.
\begin{prop} \label{henselisation-composition}
    Let $(K,K^+) = (K_1,K_1^+) \succ (K_2,K_2^+)$ be a composition of valued fields such that~$K^\succ = K_2^\succ$ is of characteristic $p \ge 0$.
    \begin{enumerate}[ref=(\arabic*)]
        \item \label{lemitem:hens-comp1}
        $(K,K^+)$ is henselian if and only if $(K_1,K_1^+)$ and $(K_2,K_2^+)$ are henselian.
        \item \label{lemitem:hens-comp2}
        $(K,K^+)$ is strictly henselian if and only if $(K_1,K_1^+)$ is henselian and $(K_2,K_2^+)$ is strictly henselian.
        
        \item \label{lemitem:hens-comp3}
        $(K,K^+)$ is tamely henselian 
        if and only if $(K_1,K_1^+)$ is henselian with value group divisible by all primes except possibly~$p$ and $(K_2,K_2^+)$ is tamely henselian.
    \end{enumerate}
\end{prop}

\begin{prop} \label{composition-tame}
    Let $(L,L^+)/(K,K^+)$ be an algebraic extension of valued fields of residue characteristic $p \ge 0$.
    Suppose that $(L,L^+) = (L_1,L_1^+) \succ (L_2,L_2^+)$ and set $(K_1,K_1^+) := (K,L_1^+ \cap K)$ and $(K_2,K_2^+) := (K_1^\succ,L_2^+ \cap K_1^\succ)$ so that $(K,K^+) = (K_1,K_1^+) \succ (K_2,K_2^\succ)$.
    Then
    \begin{enumerate}[ref=(\arabic*)]
        \item \label{lemitem:composition etale}
        $(L,L^+)/(K,K^+)$ is unramified if and only if $(L_1,L_1^+)/(K_1,K_1^+)$ and $(L_2,L_2^+)/(K_2,K_2^+)$ are unramified.
        
        In particular, the henselisation and strict henselisation of $(K,K^+)$ decompose as
        \begin{align*}
            (K^\rh,K^{\rh,+}) &= (K'_1,{K'_1}^+) \succ (K_2^\rh,K_2^{\rh,+}),   \\
            (K^\sh,K^{\sh,+}) &= (K''_1,{K''_1}^+) \succ (K_2^\sh,K_2^{\sh,+}),
        \end{align*}
        where $(K'_1,{K'_1}^+)$ and $(K''_1,{K''_1}^+)$ are the unique unramified extensions of $(K_1^\rh,K_1^{\rh,+})$ with residue fields $K_2^\rh$ and $K_2^\sh$, respectively.
        \item \label{lemitem:composition tame}
        $(L,L^+)/(K,K^+)$ is tame if and only if $(L_1,L_1^+)/(K_1,K_1^+)$ and $(L_2,L_2^+)/(K_2,K_2^+)$ are tame and the ramification index of $(L_1,L_1^+)/(K_1,K_1^+)$ is prime to~$p$ \footnote{It might happen that the residue characteristic of $(K_1,K_1^+)$ is~$0$ and the residue characteristic of $(K_2,K_2^+)$ is $p > 0$.
        In this case extracting $p$-th roots in $(K_1,K_1^+)$ is tame but does not lead to a tame extension of $(K,K^+)$.
        This is the reason for the extra condition on the ramification index.}.

        In particular, the tame henselisation of $(K,K^+)$ decomposes as
        \[
            (K^t,K^{t,+}) = (K'_1,{K'_1}^+) \succ (K_2^t,K_2^{t,+}),
        \]
         where $(K'_1,{K'_1}^+) \subseteq (K_1^t,K_1^{t,+})$ is the minimal extension of $(K_1^\rh,K_1^{\rh,+})$ with residue field~$K_2^t$ and value group divisible by all primes except possibly~$p$.
    \end{enumerate}
\end{prop}

\begin{proof}
    Let $\fp \subseteq K^+$ be the maximal ideal of~$K_1^+$, so that $K_1^+ = K^+_\fp$ and $K_2^+ = K^+/\fp$.
    
    In order to prove~\labelcref{lemitem:composition etale} suppose that $(L,L^+)/(K,K^+)$ is unramified.
    Then $\fp_L := \fp L^+$ is the maximal ideal of~$L_1^+$.
    Moreover, $L_1^+$ and $L_2^+$ are the base changes of~$L^+$ via $K^+ \hookrightarrow K_1^+$ and $K^+ \twoheadrightarrow K_2^+$, respectively.
    In particular, both extensions $L_1^+/K_1^+$ and $L_2^+/K_2^+$ are unramified.
    
    Suppose now that $(L_1,L_1^+)/(K_1,K_1^+)$ and $(L_2,L_2^+)/(K_2,K_2^+)$ are unramified.
    We can check that $(L,L^+)/(K,K^+)$ is unramified after base changing to $(K^\sh,K^{\sh,+})$.
    So we may assume that $(K,K^+)$ is strictly henselian.
    By \cref{henselisation-composition}~\labelcref{lemitem:hens-comp2}, also $(K_2,K_2^+)$ is strictly henselian and $(K_1,K_1^+)$ is henselian.
    Therefore, the unramified extension $(L_2,L_2^+)/(K_2,K_2^+)$ is trivial.
    Now $(L_1,L_1^+)/(K_1,K_1^+)$ is an unramified extension with trivial residue field extension of a henselian valued field, hence trivial as well.
    This shows that $(L,L^+) = (K,K^+)$, which is trivially unramified.

    For the second part of~\labelcref{lemitem:composition etale} we note that the unramified extension $(K^h,K^{h,+})/(K,K^+)$ induces unramified extensions of $(K_1,K_1^+)$ and $(K_2,K_2^+)$ that have to contain the respective henselisations by \cref{henselisation-composition}~\labelcref{lemitem:hens-comp1}.
    Invoking that the residue field extension over~$K_2^\succ$ has to be trivial we obtain the desired decomposition.
    The argument for the strict henselisation is analogous.

    In order to prove~\labelcref{lemitem:composition tame} we first assume that $(L,L^+)/(K,K^+)$ is tame.
    We already know from~\labelcref{lemitem:composition etale} that passing to the strict henselisation of $(K,K^+)$ induces unramified extensions of $(K_1,K_1^+)$ and $(K_2,K_2^+)$.
    We may therefore assume that $(K,K^+)$ is strictly henselian.
    Then $L/K$ is of degree prime to~$p$ and the same holds for $L_1/K_1$ and $L_2/K_2$, so both extensions $(L_1,L_1^+)/(K_1,K_1^+)$ and $(L_2,L_2^+)/(K_2,K_2^+)$ are tame and the ramification index of the former extension is prime to~$p$.
    
    For the converse we reduce to the situation where $(K,K^+)$ is tamely henselian.
    Then we argue in the same way as before using that a tame extension of henselian valued fields with trivial residue field extension and trivial value group extension is trivial.
\end{proof}

\end{document}
